% sec8_6.tex — Section 8.6: LIML

\subsection*{LIML Defined}

The advantage of the FIML estimator is that it allows you to exploit all the information afforded by the complete system of simultaneous equations. This, however,
is also a weakness because, as is true with any other system or joint estimation procedure, the estimator is not consistent if any part of the system is misspecified.
If you are confident that the equation in question is correctly specified but not so
sure about the rest of the system, you may well prefer to employ single-equation
methods such as 2SLS. The rest of this section derives the ML estimator called the
\textbf{limited-information maximum likelihood (LIML) estimator}, which is the ML
counterpart of 2SLS.

Let the $m$-th equation of the $M$-equation system be the equation of interest. The
$m$-th equation has $L_m$ regressors, $\mathbf{z}_{tm}$, are either endogenous or predetermined. Let
$\tilde{\mathbf{y}}_t$ be the vector
of $M_m$ endogenous variables and $\tilde{\mathbf{x}}_t$ be the vector of $K_m$ predetermined variables,
with $M_m + K_m = L_m$. Partition $\boldsymbol{\delta}_{0m}$ conformably as
$\boldsymbol{\delta}_{0m} = (\boldsymbol{\gamma}_{0m}', \boldsymbol{\beta}_{0m}')'$. Then the
$m$-th equation can be written as
\begin{equation}
  y_{tm} = \underset{(1 \times L_m)}{\mathbf{z}_{tm}'}
  \underset{(L_m \times 1)}{\boldsymbol{\delta}_{0m}}
  + \varepsilon_{tm}
  = \underset{(1 \times M_m)}{\tilde{\mathbf{y}}_t'}
  \underset{(M_m \times 1)}{\boldsymbol{\gamma}_{0m}}
  + \underset{(1 \times K_m)}{\tilde{\mathbf{x}}_t'}
  \underset{(K_m \times 1)}{\boldsymbol{\beta}_{0m}}
  + \varepsilon_{tm}.
  \label{eq:8.6.1} \tag{8.6.1}
\end{equation}
Obviously, this single equation in isolation is an incomplete system because of the
existence of the $M_m$ included endogenous variables $\tilde{\mathbf{y}}_t$. However, if this equation
is supplemented by the $M_m$ reduced-form equations for $\tilde{\mathbf{y}}_t$ taken from \eqref{eq:8.5.9}, then
the system of $1 + M_m$ equations thus formed is a complete system of simultaneous
equations. To describe this idea more fully, let $\tilde{\boldsymbol{\Pi}}_0$ be the associated reduced-form coefficient matrix and collect appropriate $M_m$ columns from $\boldsymbol{\Pi}_0$ and write the
reduced-form equations for the $M_m$ included endogenous variables as
\begin{equation}
  \underset{(M_m \times 1)}{\tilde{\mathbf{y}}_t}
  = \underset{(M_m \times K)}{\tilde{\boldsymbol{\Pi}}_0'}
  \underset{(K \times 1)}{\mathbf{x}_t}
  + \underset{(M_m \times 1)}{\tilde{\mathbf{v}}_t}.
  \label{eq:8.6.2} \tag{8.6.2}
\end{equation}

Combining \eqref{eq:8.6.1} and \eqref{eq:8.6.2}, we obtain the system of $1 + M_m$ equations:
\begin{equation}
  \underset{((1+M_m) \times (1+M_m))}{\bar{\boldsymbol{\Gamma}}_0}
  \underset{((1+M_m) \times 1)}{\bar{\mathbf{y}}_t}
  + \underset{((1+M_m) \times K)}{\bar{\mathbf{B}}_0}
  \underset{(K \times 1)}{\mathbf{x}_t}
  = \underset{((1+M_m) \times 1)}{\bar{\boldsymbol{\varepsilon}}_t},
  \label{eq:8.6.3} \tag{8.6.3}
\end{equation}
where
\[
  \bar{\mathbf{y}}_t \equiv \begin{bmatrix} y_{tm} \\ \tilde{\mathbf{y}}_t \end{bmatrix}, \quad
  \bar{\boldsymbol{\varepsilon}}_t \equiv \begin{bmatrix} \varepsilon_{tm} \\ \tilde{\mathbf{v}}_t \end{bmatrix},
\]
\begin{equation}
  \bar{\boldsymbol{\Gamma}}_0
  \equiv \begin{bmatrix}
    1 & -\boldsymbol{\gamma}_{0m}' \\
    \underset{(M_m \times 1)}{\mathbf{0}} & \mathbf{I}_{M_m}
  \end{bmatrix}, \quad
  \bar{\mathbf{B}}_0
  \equiv \begin{bmatrix}
    -\boldsymbol{\beta}_{0m}' & \mathbf{0}' \\
    \underset{(1 \times K_m)}{\phantom{x}} & \underset{(1 \times (K - K_m))}{\phantom{x}} \\
    & -\tilde{\boldsymbol{\Pi}}_0' \\
    \underset{(M_m \times K)}{\phantom{x}}
  \end{bmatrix}.
  \label{eq:8.6.4} \tag{8.6.4}
\end{equation}
Here, we have assumed, without loss of generality, that the included predetermined
variables $\tilde{\mathbf{x}}_t$ are the first $K_m$ elements of $\mathbf{x}_t$. The equation system \eqref{eq:8.6.3} is a complete system of $1 + M_m$ simultaneous equations because
\begin{equation}
  |\bar{\boldsymbol{\Gamma}}_0| = 1 \neq 0.
  \label{eq:8.6.5} \tag{8.6.5}
\end{equation}

For example, consider Example~8.2. The included endogenous variable in the
wage equation is $S_t$. Supplement the wage equation by a reduced-form equation
for $S_t$, a regression of $S_t$ on a constant, $IQ_t$, and $MED_t$. The $\bar{\boldsymbol{\Gamma}}_0$ matrix for the
two-equation system that results is
\[
  \bar{\boldsymbol{\Gamma}}_0
  = \begin{bmatrix} 1 & -\gamma_{11} \\ 0 & 1 \end{bmatrix}.
\]

The hypothetical value $\bar{\boldsymbol{\Gamma}}$ of $\bar{\boldsymbol{\Gamma}}_0$ has the same structure as shown in \eqref{eq:8.6.4} for
$\bar{\boldsymbol{\Gamma}}_0$. So $|\bar{\boldsymbol{\Gamma}}| = 1$. Replacing in the FIML objective function \eqref{eq:8.5.20}
$(\boldsymbol{\delta}, \mathbf{B})$
by $(\boldsymbol{\gamma}_m, \boldsymbol{\beta}_m, \bar{\boldsymbol{\Pi}})$,
$\boldsymbol{\Gamma}$ by $\bar{\boldsymbol{\Gamma}}$, and $\boldsymbol{\Sigma}$ by $\bar{\boldsymbol{\Sigma}}$, and noting that
$|\bar{\boldsymbol{\Gamma}}| = 1$, we obtain the LIML
objective function:
\begin{multline}
  Q_n(\boldsymbol{\gamma}_m, \boldsymbol{\beta}_m, \bar{\boldsymbol{\Pi}}, \bar{\boldsymbol{\Sigma}})
  = -\frac{M_m}{2}\log(2\pi)
    - \frac{1}{2}\log(|\bar{\boldsymbol{\Sigma}}|) \\
    - \frac{1}{2n}\sum_{t=1}^{n}
      [\bar{\boldsymbol{\Gamma}}\bar{\mathbf{y}}_t + \bar{\mathbf{B}}\mathbf{x}_t]'
      \bar{\boldsymbol{\Sigma}}^{-1}
      [\bar{\boldsymbol{\Gamma}}\bar{\mathbf{y}}_t + \bar{\mathbf{B}}\mathbf{x}_t],
  \label{eq:8.6.6} \tag{8.6.6}
\end{multline}
where $\bar{\boldsymbol{\Sigma}}$ is the $(1 + M_m) \times (1 + M_m)$ variance matrix of
$\bar{\boldsymbol{\varepsilon}}_t$. Let
$(\hat{\boldsymbol{\gamma}}_m, \hat{\boldsymbol{\beta}}_m, \hat{\bar{\boldsymbol{\Pi}}}, \hat{\bar{\boldsymbol{\Sigma}}})$
be the FIML estimator of this system. The \textbf{LIML estimator} of
$\boldsymbol{\delta}_{0m} = (\boldsymbol{\gamma}_{0m}', \boldsymbol{\beta}_{0m}')'$
is $(\hat{\boldsymbol{\gamma}}_m', \hat{\boldsymbol{\beta}}_m')'$. This derivation of the LIML estimator, which is different from the
original derivation of Anderson and Rubin (1949), is due to Pagan (1979). Since
LIML is a FIML estimator, Proposition~\ref{prop:8.1} assures us that the LIML estimator is
consistent and asymptotically normal even if the error is not normally distributed.

\subsection*{Computation of LIML}

For the SUR model, we have shown that the ML estimator can be obtained from
the iterative SUR procedure. If you reexamine the argument, you should see that
it does not depend on $\boldsymbol{\Gamma}$ in the SUR objective function \eqref{eq:8.5.29} being the identity
matrix; all that is needed is that the determinant of $\boldsymbol{\Gamma}$ is a constant. Therefore, the
same argument also applies to the LIML objective function \eqref{eq:8.6.6}. That is, we can
obtain the LIML estimator of $\boldsymbol{\delta}_{0m}$ by applying the iterative SUR procedure to the
$(1+M_m)$-equation system \eqref{eq:8.6.3}. In particular, if $\boldsymbol{\Sigma}_0$ is known, the LIML estimator
is the SUR estimator. On the face of it, this conclusion --- that SUR can be used to
estimate the coefficients of included \emph{endogenous} variables --- is surprising, but it is
true. The essential reason is this: the sampling error $\hat{\boldsymbol{\delta}}_m - \boldsymbol{\delta}_{0m}$ depends not only
on the correlation between the included endogenous variables $\tilde{\mathbf{y}}_t$ and the error term
$\varepsilon_{tm}$ but also on the correlation between $\tilde{\mathbf{y}}_t$ and $\tilde{\mathbf{v}}_t$. Those two correlations cancel
each other out. Review Question~1 verifies this for a simple example.

The iterative SUR procedure, however, is not how the LIML estimator is usually calculated, because there is a closed-form expression for the estimator. The
simultaneous equation system \eqref{eq:8.6.3} has two special features. One is the special
structure of $\bar{\boldsymbol{\Gamma}}$, which we have noted, and the other is the fact that there are no
exclusion restrictions in the rows of $\bar{\mathbf{B}}$ corresponding to the included endogenous
variables $\tilde{\mathbf{y}}_t$. Thanks to these features, the LIML estimator
$\hat{\boldsymbol{\delta}}_m = (\hat{\boldsymbol{\gamma}}_m', \hat{\boldsymbol{\beta}}_m')'$ and
also the likelihood ratio statistic \eqref{eq:8.5.25} can be calculated explicitly.

To write down the closed-form expressions for the LIML estimator and the
likelihood ratio statistic, we need to introduce some more notation. Let $\mathbf{Z}_m$,
$\tilde{\mathbf{Y}}$, $\tilde{\mathbf{X}}$,
$\mathbf{X}$, and $\mathbf{y}_m$ be the data matrices and the data vector associated with $\mathbf{z}_{tm}$,
$\tilde{\mathbf{y}}_t$, $\tilde{\mathbf{x}}_t$, $\mathbf{x}_t$,
and $y_{tm}$, respectively:
\begin{equation}
  \underset{(n \times L_m)}{\mathbf{Z}_m}
  \equiv \begin{bmatrix} \mathbf{z}_{1m}' \\ \vdots \\ \mathbf{z}_{nm}' \end{bmatrix}, \quad
  \underset{(n \times M_m)}{\tilde{\mathbf{Y}}}
  \equiv \begin{bmatrix} \tilde{\mathbf{y}}_1' \\ \vdots \\ \tilde{\mathbf{y}}_n' \end{bmatrix}, \quad
  \underset{(n \times K_m)}{\tilde{\mathbf{X}}}
  \equiv \begin{bmatrix} \tilde{\mathbf{x}}_1' \\ \vdots \\ \tilde{\mathbf{x}}_n' \end{bmatrix}, \quad
  \underset{(n \times K)}{\mathbf{X}}
  \equiv \begin{bmatrix} \mathbf{x}_1' \\ \vdots \\ \mathbf{x}_n' \end{bmatrix}, \quad
  \underset{(n \times 1)}{\mathbf{y}_m}
  \equiv \begin{bmatrix} y_{1m} \\ \vdots \\ y_{nm} \end{bmatrix},
  \label{eq:8.6.7} \tag{8.6.7}
\end{equation}
and let $\mathbf{M}$ and $\tilde{\mathbf{M}}$ be the annihilators based on $\mathbf{X}$ and $\tilde{\mathbf{X}}$, respectively:
\begin{equation}
  \underset{(n \times n)}{\mathbf{M}}
  \equiv \mathbf{I}_n - \mathbf{X}(\mathbf{X}'\mathbf{X})^{-1}\mathbf{X}', \quad
  \underset{(n \times n)}{\tilde{\mathbf{M}}}
  \equiv \mathbf{I}_n - \tilde{\mathbf{X}}(\tilde{\mathbf{X}}'\tilde{\mathbf{X}})^{-1}\tilde{\mathbf{X}}'.
  \label{eq:8.6.8} \tag{8.6.8}
\end{equation}
Then the LIML estimator $\hat{\boldsymbol{\delta}}_m$ can be written as
\begin{equation}
  \hat{\boldsymbol{\delta}}_m
  = [\mathbf{Z}_m'(\mathbf{I}_n - k\mathbf{M})\mathbf{Z}_m]^{-1}
    \mathbf{Z}_m'(\mathbf{I}_n - k\mathbf{M})\mathbf{y}_m,
  \label{eq:8.6.9} \tag{8.6.9}
\end{equation}
where $k$ is the smallest characteristic root of
\begin{equation}
  \tilde{\mathbf{W}}\,\mathbf{W}^{-1},
  \label{eq:8.6.10} \tag{8.6.10}
\end{equation}
with $(1 + M_m) \times (1 + M_m)$ matrices $\tilde{\mathbf{W}}$ and $\mathbf{W}$ defined by
\begin{equation}
  \tilde{\mathbf{W}}
  \equiv [\mathbf{y}_m : \tilde{\mathbf{Y}}]'\tilde{\mathbf{M}}
    [\mathbf{y}_m : \tilde{\mathbf{Y}}], \quad
  \mathbf{W}
  \equiv [\mathbf{y}_m : \tilde{\mathbf{Y}}]'\mathbf{M}
    [\mathbf{y}_m : \tilde{\mathbf{Y}}].
  \label{eq:8.6.11} \tag{8.6.11}
\end{equation}
(See, e.g., Davidson and MacKinnon, 1993, pp.~645--649, for a derivation.) The
likelihood ratio statistic \eqref{eq:8.5.25} for testing overidentifying restrictions reduces to
\begin{equation}
  LR = n \log k.
  \label{eq:8.6.12} \tag{8.6.12}
\end{equation}
Since there are no overidentifying restrictions in the supplementary $M_m$ equations
\eqref{eq:8.6.2}, the equation of interest is the only source of overidentifying restrictions,
which are $K - L_m$ in number. Therefore, by Proposition~\ref{prop:8.1}, this statistic for testing
overidentifying restrictions is asymptotically $\chi^2(K - L_m)$. When the equation is
just identified so that $K = L_m$, then the $LR$ statistic should be zero. Indeed, it can
be shown that $k = 1$ in this case (see, e.g., p.~647 of Davidson and MacKinnon,
1993).

If we do not necessarily require $k$ to be as just defined, the estimator \eqref{eq:8.6.9} is
called a \textbf{$k$-class estimator}. The LIML estimator is thus a $k$-class estimator with
$k$ defined above. Inspection of the 2SLS formula in terms of data matrices (see
$(3.8.3')$ on page~230) immediately shows that the 2SLS estimator is a $k$-class estimator with $k = 1$, and the OLS estimator is a $k$-class estimator with $k = 0$. It
follows that LIML and 2SLS are numerically the same when the equation is just
identified (so that $k = 1$).

We have already shown that LIML is consistent and asymptotically normal.
Moreover, using the explicit formula \eqref{eq:8.6.9}, it can be shown that the LIML estimator of $\boldsymbol{\delta}_{0m}$ has the same asymptotic distribution as 2SLS (see, e.g., Amemiya,
1985, Section~7.3.4). So the formula for the asymptotic variance of LIML is given
by (3.8.4) and its consistent estimate by (3.8.5).

\subsection*{LIML versus 2SLS}

Since LIML and 2SLS share the same asymptotic distribution, you cannot prefer
one over the other on asymptotic grounds. For any finite sample, LIML has the
invariance property, while 2SLS is not invariant. Furthermore, the conclusion one
could draw from the large literature on finite-sample properties (see, e.g., Judge
et al., 1985, Section~15.4) is that LIML should be preferred to 2SLS. Major conclusions from the literature are listed below. They assume that the predetermined
variables $\mathbf{x}_t$ are fixed constants, however.
\begin{itemize}
\item Suppose errors are normally distributed. The $p$-th moment of the 2SLS estimator exists if and only if $p < K - L_m + 1$, where $K$ is the number of predetermined
  variables and $L_m$ is the number of regressors in the $m$-th equation. Thus, 2SLS
  does not even have a mean if the equation is just identified (i.e., if $K = L_m$),
  and this is true even if the distribution of errors is a normal distribution, whose
  moments are all finite.

\item The LIML estimator has no finite moments even if errors are normally distributed.

\item However, in most Monte Carlo simulations (see, e.g., Anderson, Kunitomo, and
  Sawa, 1982), LIML approaches its asymptotic normal distribution much more
  rapidly than does 2SLS.
\end{itemize}

\bigskip
\noindent\rule{\textwidth}{0.4pt}
\textbf{QUESTIONS FOR REVIEW}
\medskip

\begin{enumerate}
\item (Why SUR can be used for a structural equation)\quad
  To see why SUR can be
  used to estimate the coefficients of included endogenous variables, consider
  the following simple example:
  \[
    \begin{cases}
      y_{t1} = \gamma_0 y_{t2} + \varepsilon_{t1}, \\
      y_{t2} = \beta_0 x_t + v_t,
    \end{cases}
  \]
  where $x_t$ is predetermined. Suppose for simplicity that
  \[
    \underset{(2 \times 2)}{\boldsymbol{\Sigma}_0}
    \equiv \E\!\left(\begin{bmatrix} \varepsilon_{t1} \\ v_t \end{bmatrix}
      \begin{bmatrix} \varepsilon_{t1} & v_t \end{bmatrix}\right)
  \]
  is known. Let $(\hat{\gamma}, \hat{\beta})$ be the SUR estimator of $(\gamma_0, \beta_0)$.
\end{enumerate}
